\chapter{Discuss\~ao, Limita\c{c}\~oes e Conclus\~ao}
\label{cap:conclusao}

Este cap\'itulo retoma os principais resultados do trabalho, explicita suas limita\c{c}\~oes, sintetiza a contribui\c{c}\~ao alcan\c{c}ada e apresenta possibilidades de continuidade. O objetivo \'e fechar a narrativa do TCC de forma coerente com a evolu\c{c}\~ao do projeto e com o uso institucional pretendido para a solu\c{c}\~ao.

\section{Discuss\~ao dos resultados}
\label{sec:discussao}
\label{sec:retomada-problema}
\label{sec:atendimento-objetivos}
\label{sec:principais-contribuicoes}
\label{sec:aspectos-eticos-uso-responsavel}

Os resultados mostram que a solu\c{c}\~ao final apresentou ader\^encia melhor ao problema do que as hip\'oteses iniciais do projeto. Embora grafos e LSTM tenham sido importantes na fase explorat\'oria, a experi\^encia emp\'irica indicou que o problema central estava melhor representado por uma base tabular temporal com atributos derivados e valida\c{c}\~ao cronol\'ogica. Essa conclus\~ao \'e relevante porque n\~ao foi assumida de antem\~ao: ela surgiu do confronto entre diferentes caminhos metodol\'ogicos e os limites concretos da base escolar analisada.

Outro ponto importante \'e o valor dos baselines. Em diferentes momentos do trabalho, regras simples baseadas na \'ultima nota ou em m\'edias m\'oveis competiram de forma relevante com modelos mais sofisticados. Isso refor\c{c}a que complexidade algor\'itmica n\~ao deve ser tomada como sin\^onimo autom\'atico de superioridade. No contexto deste TCC, esse cuidado ajudou a manter a an\'alise mais honesta e mais pr\'oxima do uso real da solu\c{c}\~ao.

Do ponto de vista institucional, o maior avan\c{c}o do projeto n\~ao est\'a apenas nas m\'etricas, mas na passagem da predi\c{c}\~ao para a a\c{c}\~ao. O valor do sistema aparece quando a nota prevista, o alerta de risco e os sinais complementares s\~ao convertidos em relat\'orios intelig\'iveis para professor, coordena\c{c}\~ao e secretaria. Na pr\'atica, esses alertas devem orientar investiga\c{c}\~ao pedag\'ogica, prioriza\c{c}\~ao de contato, planejamento de refor\c{c}o e acompanhamento da evolu\c{c}\~ao do aluno. Em outras palavras, a contribui\c{c}\~ao do trabalho se sustenta mais quando a modelagem \'e lida como apoio \`a decis\~ao do que como mecanismo aut\^onomo de julgamento.

\section{Limita\c{c}\~oes do trabalho}
\label{sec:limitacoes}

As principais limita\c{c}\~oes identificadas foram:
\begin{enumerate}
    \item o v\'inculo de professor ainda n\~ao cobre integralmente todo o hist\'orico analisado;
    \item o ranking executivo ainda concentra muitos casos nas faixas mais altas de prioridade;
    \item a base de 2026 j\'a existe, mas nem todos os cen\'arios possuem alvo consolidado suficiente para teste completo;
    \item o projeto utiliza vari\'aveis institucionais dispon\'iveis, mas n\~ao incorpora dimens\~oes qualitativas, socioemocionais ou contextuais mais profundas;
    \item a generaliza\c{c}\~ao para outras escolas exige adapta\c{c}\~ao local da arquitetura e nova valida\c{c}\~ao.
\end{enumerate}

Tamb\'em existem amea\c{c}as \`a validade interna, externa e de constru\c{c}\~ao. Nem toda queda de desempenho \'e explicada pelas vari\'aveis dispon\'iveis, e conceitos educacionais complexos acabam representados por \textit{proxies} quantitativos, como faltas, pagamentos e trajet\'oria de notas. Essa limita\c{c}\~ao n\~ao invalida o projeto, mas delimita com mais clareza o alcance de suas interpreta\c{c}\~oes.

\section{Conclus\~ao}
\label{sec:conclusao-final}

Este Trabalho de Conclus\~ao de Curso apresentou o desenvolvimento de um sistema preditivo para apoio \`a decis\~ao pedag\'ogica com base em dados escolares reais. O projeto evoluiu de uma hip\'otese inicial orientada a grafos para uma pipeline temporal tabular reproduz\'ivel, capaz de:
\begin{itemize}
    \item prever a pr\'oxima nota do aluno com desempenho superior a baselines simples;
    \item identificar risco de nota futura abaixo de 6 com \textit{recall} elevado;
    \item integrar notas, faltas, pagamentos e contexto escolar em uma arquitetura \'unica;
    \item transformar resultados t\'ecnicos em relat\'orios operacionais para diferentes atores da escola.
\end{itemize}

Em termos quantitativos, o modo de previs\~ao num\'erica atingiu MAE de 0,7659, enquanto o modo de alerta de risco alcan\c{c}ou F1 de 0,7436 e \textit{recall} de 0,8794 no teste temporal. Em termos pr\'aticos, esses resultados indicam que a solu\c{c}\~ao j\'a pode servir como mecanismo de triagem e prioriza\c{c}\~ao institucional, desde que usada com supervis\~ao humana e interpreta\c{c}\~ao contextualizada. O ponto central n\~ao \'e afirmar que a escola passa a ``saber o futuro'' do aluno, mas que passa a dispor de um instrumento adicional para agir com mais anteced\^encia.

Assim, a principal contribui\c{c}\~ao do trabalho n\~ao est\'a apenas em acertar melhor uma nota futura, mas em mostrar que \'e poss\'ivel construir uma arquitetura anal\'itica coerente com o cotidiano da escola, eticamente defens\'avel e operacionalmente \'util. Nesse sentido, o TCC procura contribuir menos com uma promessa de solu\c{c}\~ao definitiva e mais com uma forma consistente de integrar dados, modelagem e uso institucional em educa\c{c}\~ao.

\section{Trabalhos futuros}
\label{sec:futuro}

Como continuidade natural do projeto, o principal desdobramento previsto \'e consolidar o sistema como uma plataforma de relat\'orios operacionais para col\'egios e faculdades. A proposta \'e evoluir os artefatos demonstrados neste TCC para um fluxo recorrente de acompanhamento, no qual a modelagem alimente relat\'orios voltados a decis\~oes pedag\'ogicas e administrativas.

Nesse cen\'ario, coordenadores, professores e secretaria passariam a consultar relat\'orios espec\'ificos para identificar alunos com maior risco de queda, disciplinas com concentra\c{c}\~ao de alertas, turmas que exigem interven\c{c}\~ao coletiva e casos que demandam contato administrativo. O objetivo \'e antecipar medidas de apoio antes que o problema se consolide como reprova\c{c}\~ao bimestral, trimestral, semestral ou anual.

Para isso, os relat\'orios atuais devem evoluir de sa\'idas anal\'iticas do TCC para um mecanismo institucional mais est\'avel, com ciclos de atualiza\c{c}\~ao definidos, filtros por unidade, curso, turma e disciplina, e trilhas de acompanhamento que permitam verificar se o caso melhorou ap\'os a interven\c{c}\~ao. Em faculdades, a mesma arquitetura pode ser adaptada a componentes curriculares semestrais, disciplinas com alta reprova\c{c}\~ao hist\'orica, sinais de evas\~ao e combina\c{c}\~oes entre desempenho acad\^emico e informa\c{c}\~oes administrativas.

Outra amplia\c{c}\~ao relevante ser\'a obter dados acad\^emicos mais detalhados sobre tarefas e atividades realizadas pelos alunos. Al\'em das notas formais, faltas e pagamentos, futuras vers\~oes podem incorporar informa\c{c}\~oes sobre tarefas propostas, tarefas realizadas, entregas no prazo, entregas atrasadas, atividades n\~ao entregues e notas ou conceitos atribu\'idos a essas atividades. O cruzamento desses sinais com as m\'edias escolares e com o hist\'orico de desempenho pode revelar ind\'icios mais precoces de queda, dificuldade de organiza\c{c}\~ao, desengajamento ou lacunas de aprendizagem, fortalecendo a capacidade do sistema de antecipar risco acad\^emico antes da reprova\c{c}\~ao.

Tamb\'em ser\'a necess\'ario aprofundar a escolha do algoritmo mais adequado quando o sistema for testado em diferentes col\'egios e faculdades. Embora este trabalho tenha consolidado uma pipeline tabular reproduz\'ivel, a decis\~ao final entre modelos como \textit{Random Forest} e \textit{HistGradientBoosting} deve considerar o comportamento observado em bases institucionais distintas, com diferentes volumes de dados, padr\~oes de avalia\c{c}\~ao, distribui\c{c}\~ao de faltas, organiza\c{c}\~ao de turmas e hist\'oricos de reprova\c{c}\~ao. Essa compara\c{c}\~ao deve manter o mesmo cuidado adotado neste TCC: valida\c{c}\~ao temporal, compara\c{c}\~ao com baselines simples e an\'alise do custo pedag\'ogico dos erros.

Por fim, a evolu\c{c}\~ao para uso operacional exigir\'a ajustes no sistema RP da institui\c{c}\~ao, especialmente na defini\c{c}\~ao de regras, telas e notifica\c{c}\~oes que transformem o risco previsto em a\c{c}\~ao concreta. Um exemplo previsto \'e a cria\c{c}\~ao de alertas no aplicativo m\'ovel do professor, permitindo que ele seja avisado quando um aluno apresentar risco acad\^emico relevante e possa consultar os principais sinais associados ao alerta. Esse recurso deve ser implementado com cuidado para evitar excesso de notifica\c{c}\~oes, preservar a interpreta\c{c}\~ao pedag\'ogica e manter a decis\~ao final sob responsabilidade humana.

Al\'em disso, permanecem relevantes os seguintes desdobramentos t\'ecnicos:
\begin{enumerate}
    \item ampliar a cobertura do v\'inculo com professores e reduzir casos sem identifica\c{c}\~ao docente;
    \item recalibrar os limiares do ranking executivo e da pontua\c{c}\~ao composta para diminuir sobrecarga operacional;
    \item incorporar t\'ecnicas adicionais de explicabilidade para justificar melhor cada alerta;
    \item medir o impacto real dos relat\'orios nas decis\~oes de coordenadores, professores e secretaria;
    \item comparar \textit{Random Forest} e \textit{HistGradientBoosting} em bases de diferentes institui\c{c}\~oes para definir qual modelo se ajusta melhor ao uso pretendido;
    \item incorporar dados de tarefas, entregas, atrasos, n\~ao entregas e notas de atividades, cruzando esses sinais com as m\'edias escolares e com o hist\'orico de desempenho do aluno;
    \item integrar os alertas ao sistema RP e ao aplicativo m\'ovel do professor, com regras de notifica\c{c}\~ao, consulta dos motivos do alerta e controle de sobrecarga operacional;
    \item validar a solu\c{c}\~ao em outras escolas, cursos superiores ou redes educacionais.
\end{enumerate}
